\uuid{4rCp}
\titre{Suite arithmétique/géométrique}
\auteur{bodin ; hanani ; mzari}
\organisation{exo7}
\licence{CC BY-NC-SA 4.0}
\datecreate{2020-01-20}
\niveau{L1}
\module{Analyse}
\chapitre{Suite}
\sousChapitre{Autre}
\difficulte{1}
\theme{Suites, suite arithmétique, géométrique}
\exo7qcmid{lille-1.suites.008}

\contenu{
\begin{question}
\qclassification{121.00}
\qtags{motcle=suite arithmétique/géométrique}
Soit $\displaystyle u_n=1+\frac{1}{2}+\frac{1}{2^2}+\dots +\frac{1}{2^n}$. Quelles sont les bonnes réponses ?
\begin{answers}
    \bad{La suite $(u_n)$ est divergente.}
    \good{La suite $(u_n)$ est strictement croissante.}
    \bad{$\displaystyle \lim _{n\to +\infty}u_n=+\infty$}
    \good{$\displaystyle \lim _{n\to +\infty}u_n=2$}
\end{answers}
\begin{explanations}
Le terme $u_n$ est la somme des premiers termes de la suite géométrique de raison $\displaystyle \frac{1}{2}$. Donc $(u_n)$ est strictement croissante et
$$u_n=\frac{1-\frac{1}{2^{n+1}}}{1-\frac{1}{2}}=2-\frac{1}{2^n}\underset{+\infty}{\longrightarrow }2.$$
\end{explanations}
\end{question}
}
